\section{异构无人集群的研究对象与问题分类}
\label{sec:01_06}

\subsection{设备类型、能力与角色资格}

\textbf{行动主体与能力差异。}
巡检系统中的无人机承担空中观测，车辆承担道路运输与地面保障，相机提供观测模态，电池提供具有身份和周转状态的能源。车辆与无人机是行动主体，相机与电池则作为载荷或资源进入决策，不能把所有设备数量都解释为机器人数量。设备异构可体现为运动方式、传感模态、服务接口及性能边界不同；同型设备也可能因位置、电量和故障状态不同而具有不同的当前可用性。前者描述配置，后者描述状态，两者都影响分配，但所需证据不同。

\textbf{资格与可执行性。}
先用资格矩阵判断设备能否承担某类角色，再检查实际位置、视线、姿态、能量和时间条件。STRATA区分可累加与不可累加特征，并将不可累加特征按最低要求转为资格指标；指标之和统计达标个体数量，不表示速度或专门技能本身相加\cite{STRATA2020}。本节采用确定性二元资格说明这一关系，不引入其随机特征分布。

设参与某项任务的机器人集合为$\mathcal C$，机器人$a$对角色$k$的资格为$q_{ak}\in\{0,1\}$，均无量纲。当角色只要求至少一个成员具备该项能力时，其必要条件为
\begin{equation}
\max_{a\in\mathcal C}q_{ak}=1,
\qquad \mathcal C\ne\varnothing.
\label{s0106:role_presence}
\end{equation}
这不保证该成员能够到达作业位置；若多个角色必须由同一设备承担，还须对同一成员同时检查各项资格，不能分别寻找不同成员后直接宣布可行。

\textbf{教学例A：角色资格下的完整分配。}
本书构造三台机器人：可见光无人机$U_1$、可见光兼热成像无人机$U_2$和地面保障车辆$V_1$。三个任务依次为可见光巡检、热成像巡检、地面保障。全部资格输入见表\ref{s0106:eligibility}；数值仅为教学设定，不是区域设备实测参数。
\begin{table}[htbp]
\centering
\caption{教学例A的能力资格矩阵（无量纲）}
\label{s0106:eligibility}
\begin{tabular}{lccc}
\toprule
设备 & 可见光 & 热成像 & 地面保障 \\
\midrule
$U_1$ & 1 & 0 & 0 \\
$U_2$ & 1 & 1 & 0 \\
$V_1$ & 0 & 0 & 1 \\
\bottomrule
\end{tabular}
\end{table}

按表中顺序，将机器人和任务分别编号为$1,2,3$。令$\alpha_{ai}=1$表示机器人$a$承担任务$i$，否则为零。本例要求每个任务恰由一台机器人承担，每台机器人恰承担一个任务；不设置收益、服务时长、路径或电池参数，求解目标仅为寻找全部可行分配。完整数值模型为
\begin{align}
\alpha_{11}+\alpha_{12}+\alpha_{13}&=1, &
\alpha_{11}+\alpha_{21}+\alpha_{31}&=1,\notag\\
\alpha_{21}+\alpha_{22}+\alpha_{23}&=1, &
\alpha_{12}+\alpha_{22}+\alpha_{32}&=1,\notag\\
\alpha_{31}+\alpha_{32}+\alpha_{33}&=1, &
\alpha_{13}+\alpha_{23}+\alpha_{33}&=1,\notag\\
\alpha_{12}=\alpha_{13}=\alpha_{23}
=\alpha_{31}=\alpha_{32}&=0,\notag\\
\alpha_{ai}&\in\{0,1\}, & a,i&\in\{1,2,3\}.
\label{s0106:example_model}
\end{align}
左右两列等式分别限制机器人负荷和任务覆盖，零值约束排除不合资格的角色。保存的Python穷举记录检查了全部六个排列，结果见表\ref{s0106:enumeration}。表中三元组按任务顺序记录承担它们的机器人编号。
\begin{table}[htbp]
\centering
\caption{教学例A的全部排列及资格检查}
\label{s0106:enumeration}
\begin{tabular}{ccc}
\toprule
分配排列 & 可行 & 不合资格的任务编号 \\
\midrule
$(1,2,3)$ & 是 & --- \\
$(1,3,2)$ & 否 & $2,3$ \\
$(2,1,3)$ & 否 & $2$ \\
$(2,3,1)$ & 否 & $2,3$ \\
$(3,1,2)$ & 否 & $1,2,3$ \\
$(3,2,1)$ & 否 & $1,3$ \\
\bottomrule
\end{tabular}
\end{table}

唯一可行分配为$U_1$承担可见光巡检、$U_2$承担热成像巡检、$V_1$承担地面保障。虽然$U_2$具备两种成像资格，本例的单任务限制仍禁止它同时占据两个任务位置。可见光资格数量不能替代热成像资格，资格数量也不能替代起降位容量。一般资源约束仍须逐资源写为$\sum_{e\in E}A_{ve}x_e\le b_v$，其中$A_{ve}$表示候选是否占用资源$v$，$b_v$为该资源允许的占用数；不同地点、接口或身份的资源不能未经可替代性证明就合并容量。本例只是有限资格分配的穷举核验，不是商用求解器实验或公开算法复现。

\subsection{团队组织与群集行为}

\textbf{团队与群集。}
本书用团队强调共同目标下的角色分工和协作关系，用群集讨论大量个体通过局部交互形成整体行为的组织方式。二者不是按固定数量阈值划分的互斥类别，也不分别等同于异构和同构。设备不同并不自动产生分布式控制；存在局部控制器，也不说明任务分配由局部交互产生。

因此，描述系统还需交代谁掌握任务和状态、谁作出分配、谁执行动作，以及信息怎样返回决策端。当前车机巡检在联合优化层面统一安排分配、航次和资源周转，应按具有明确设备身份的协作团队理解。STRATA以种群特征和任务间转移率调节群体分布\cite{STRATA2020}，体现了另一种组织粒度；其群体分布满足能力需求，不直接给出本书所需的逐机连续航迹与逐电池周转记录。

\subsection{任务依赖与MRTA分类}

\textbf{任务依赖。}
任务依赖至少包括先后关系、共同资源限制和共同完成条件。先后关系要求后续动作等待前序服务；资源限制使两个各自可行的动作不能同时执行；共同完成条件则使局部完成不能直接计作整体完成。巡检中，起降与车辆停靠相连，下一航次与飞机及电池释放相连，整塔完成与塔内全部必要观测相连。遗漏这些关系会改变可行域或完成收益。iTax将相互关联的效用与约束纳入任务分配分类\cite{Korsah2013}。

\textbf{任务分配分类。}
多机器人任务分配（MRTA）可按单机器人任务与多机器人任务（SR/MR）、单任务机器人与多任务机器人（ST/MT）、即时分配与时间扩展分配（IA/TA）区分\cite{STRATA2020}。SR/MR询问一个任务是否需要多个机器人协作；ST/MT询问一个机器人能否同时执行多个任务；IA/TA询问分配是否显式考虑跨时段安排。机器人在一天内依次执行多个任务，并不因此成为MT。

教学例A属于ST--SR--IA：每台设备只承担一个任务，每项任务只需一台设备，模型不含时序。当前巡检若以单次观测为任务，可按单机观测、时间扩展分配理解；若将车机共同完成的起降与保障活动合为服务单元，则须表达多主体配合。分类依赖任务边界，不能仅因系统中同时有车和飞机，就把每次拍摄都称为多机器人任务。进一步把分配与连续运动相连，符合任务与运动规划关注的离散决策和连续可行性结合问题\cite{TAMP2021}；这里的联系不构成对某个规划算法的性能结论。

\subsection{当前巡检研究的位置与证据范围}

\textbf{巡检研究范围。}
参考场景包含237座塔、2133项观测任务、2辆车、4架无人机和8对物理电池，班次长度为$28800\,\mathrm{s}$。地理位置取自保留的OpenStreetMap数据\cite{OSMData}，观测目标几何和部分设备参数采用模板或明示参考值。当前已知异构性主要来自地面车辆与空中无人机；无人机内部采用同一参考配置。教学例A中的传感模态差异不能据此转写为区域中已实现的不同机型配置。

研究对象是允许部分完成的车机巡检联合决策，目标依次比较未完成塔数、全队返回用时和机载能耗。已有组件验证与区域有限转录的历史记录，当前参考物理模型尚未形成新的237塔完整联合求解与执行闭环。教学分配的穷举结果仅验证本例的资格约束，不能代替完整区域中的运动、能源和协作可行性。
